\section{The certified theta moments as observations of the total original object}
\label{GraphV2:sec:certified-hankel-attachment}

The two certified matrices attach to the same original theta seed
and the same finite arithmetic packets in $\mathbf C_U$. The
global moment sum includes critical and offcritical packets in the
ambient $Q=\mathscr B/\Theta V$ of (TO.5). Its offcritical part is
exactly the already constructed $T_UQ$ of (TO.7), not a replacement
for that carrier. We construct the attachment and calculate what
the finite certificate excludes.

\subsection{From the original theta boundary to the finite moment matrices}

Retain the original operators and seed, with every scalar unchanged:
\[
 \begin{gathered}
 D=-x\partial_x,\quad
 \phi_*=(4\pi^2x^4-6\pi x^2)e^{-\pi x^2}\in V,\quad
 f_0=\Theta\phi_*=2\sum_{n\geq1}\phi_*(nx),\\
 \mathcal MF(s)=\int_0^\infty F(x)x^s\frac{dx}{x},\quad
 \mathcal Mf_0=g=2\xi,\qquad
 \psi(y)=e^{y/2}f_0(e^y).
 \end{gathered}
 \tag{HCA.1}
\]
The complete theta calculation (WBR1)--(WBR8) gives
$\psi(-y)=\psi(y)>0$ and, for $y\geq0$,
$\psi(y)\leq C_\theta e^{9y/2-\pi e^{2y}}$ with
$C_\theta=8\pi^2\sum_{n\geq1}n^4e^{-\pi(n^2-1)}<\infty$.
It follows by differentiation under this integrable majorant that
\[
 M_{2j}=\int_{\mathbb{R}}y^{2j}\psi(y)\,dy=g^{(2j)}(1/2),
 \quad a_j=\frac{(-1)^jM_{2j}}{(2j)!},\quad
 F(w)=\sum_{j\geq0}a_jw^j,\quad F(z^2)=g(1/2+iz).
 \tag{HCA.2}
\]
The map $z\mapsto z^2$ and the even entire Taylor series define
$F$ without choosing a square-root branch. Its mass is exactly
$a_0=M_0=g(1/2)>0$.

The original moment observation on a general $F_1\in\mathscr B$ is
the linear map
\[
 F_1\longmapsto
 \left(\int_{\mathbb{R}}y^{2j}e^{y/2}F_1(e^y)\,dy\right)_{j=0}^{32}.
 \tag{HCA.3}
\]
The two-ended defining bounds of $\mathscr B$ dominate each
integrand by an integrable exponential times a polynomial, so every
component is defined. At the actual boundary $f_0=\Theta\phi_*$,
this map has first component $M_0>0$. Thus it is an observation of
the actual degree-one source functions, and does not factor through
their theta quotient: factoring would send the class of $f_0$ to zero.
This explicit map is permitted by the marked source vertices of
(TO.1)--(TO.4), which retain $\phi_*$, $f_0$ and their differential.

On the open set of coefficient vectors with $a_0\ne0$, define the
finite rational recurrence
\[
 b_n=\frac{-(n+1)a_{n+1}-\sum_{j=1}^na_jb_{n-j}}{a_0}
 \quad(0\leq n\leq31),\qquad
 H_{15}^{(\nu)}=(b_{r+s+\nu})_{0\leq r,s\leq15},\quad\nu=0,1.
 \tag{HCA.4}
\]
Multiplication of $-F'/F=\sum_{n\geq0}b_nw^n$ by $F$ proves
each coefficient formula. Thus the source observation (HCA.3),
the factorial/sign map (HCA.2), the recurrence and the two matrix
assemblies form a single explicit composite. The recurrence is a
rational observation of its coefficients; it is not asserted to
be a linear or cochain map. This is exactly the typed convention
for numerical observations in (TO.4)--(TO.6).

\subsection{The reciprocal action on every full original packet}

Let $Z$ be any finite actual zero packet in the open strip, stable
under reflection and conjugation, with all full orders $m_\rho$.
Critical points are allowed in this ambient packet. Put
\[
 \begin{gathered}
 h(s)=\prod_{\rho\in Z}(s-\rho)^{m_\rho},\quad
 E_h=\mathbb{C}[s]/(h),\quad A_h=M_s,\quad v_h=g/h,\\
 \upsilon_h=j_h(v_h),\quad\varepsilon_h=\upsilon_h^{-1},\quad
 R_h(u)=r_h(\varepsilon_hu),\\
 s_hu=S_h\Theta(R_h(u)(D)\phi_*),\quad
 \sigma_hu=[s_hu]\in Q,\quad J_h=j_h\mathcal M.
 \end{gathered}
 \tag{HCA.5}
\]
Here $r_h$ is the unique degree-below-$\deg h$ remainder and $S_h$
is the original iterated Euler inverse on zero full Mellin jets.
These are precisely (AG6)--(AG8) and (WBR22)--(WBR25): the
workbench's $\eta_h$ is the same element as $\varepsilon_h$ here.
In particular
\[
 \mathcal Ms_hu=v_hR_h(u),\quad J_hs_h=I,
 \quad J_h\Theta=0,\quad D_Q\sigma_h=\sigma_hA_h.
 \tag{HCA.6}
\]
The first equality follows by dividing $gR_h(u)$ by $h$; the
second is $j_h(v_h)\varepsilon_h=1$ with the entire Taylor unit
retained. Thus $\sigma_h$ is injective. Polynomial division gives
$Ds_h-s_hA_h=\Theta\phi_*[s^{\deg h-1}]R_h$, proving the
last equality before passing to the quotient, including its boundary.
The range and inverse proofs used here apply to every root in
the open strip and therefore include the ambient critical packets.

Since $g(1/2)>0$, the polynomial $s-1/2$ is a unit on every
$E_h$. Define the actual finite-packet operator
\[
 T_h=-(A_h-\tfrac12I)^{-2},\qquad
 \beta_\rho=-\frac1{(\rho-1/2)^2}.
 \tag{HCA.7}
\]
No inverse on the whole quotient $Q$ is required. On the full
$\rho$-block, write $A_h=\rho I+N_\rho$ and
$\lambda_\rho=\rho-1/2$. The finite coordinate maps are
\[
 \begin{gathered}
 T_h=\sum_{j=0}^{m_\rho-1}(-1)^{j+1}(j+1)
                 \lambda_\rho^{-j-2}N_\rho^j,
 \quad Y_\rho=T_h-\beta_\rho I,\\
 N_\rho=\lambda_\rho\sum_{j=1}^{m_\rho-1}
             \binom{-1/2}{j}(Y_\rho/\beta_\rho)^j.
 \end{gathered}
 \tag{HCA.8}
\]
Multiplication of the first polynomial by
$(\lambda_\rho+N_\rho)^2$ gives $-1$ modulo $N_\rho^{m_\rho}$.
For the inverse, the element with constant term one solving
$U^2(1+Y_\rho/\beta_\rho)=1$ is unique by successive coefficient
comparison, whose leading scalar is two. The binomial polynomial
and $1+N_\rho/\lambda_\rho$ both solve it. This proves (HCA.8)
with every nilpotent coefficient and with no analytic branch choice.

For $h\mid H$ with full orders, let $\iota_{hH}$ insert zero
jets at added centres, let $\pi_{Hh}$ restrict to old jets, and
let $e_{hH}=\iota_{hH}(1)$ be the complete old-packet idempotent.
The original arrows (AG10)--(AG12) give
\[
 \begin{gathered}
 \pi\iota=I,\quad\iota\pi=M_{e_{hH}},\quad
 s_H\iota=s_h,\quad\sigma_H\iota=\sigma_h,\\
 \pi T_H=T_h\pi,\qquad T_H\iota=\iota T_h.
 \end{gathered}
 \tag{HCA.9}
\]
To prove the last two identities, the maps commute with $A-1/2$;
its invertibility on both domains makes them commute with its
inverse, and then with its inverse square. Thus the reciprocal
action is natural for the actual divisor arrows, including their
full units and literal representatives. Through $\sigma_h$ it
acts on the actual finite image in $Q$.

The representative of each polynomial of this action is also fixed.
Let $t_h(s)=r_h(-(s-1/2)^{-2})$, let
$p_h(s)=r_h(P(t_h(s)))$, and put
\[
 K_{P,h}(u;s)=
 \frac{p_h(s)R_h(u)(s)-R_h(P(T_h)u)(s)}{h(s)}.
 \tag{HCA.10}
\]
The numerator is zero modulo $h$ because multiplication by the
unit $\varepsilon_h$ commutes with $P(T_h)$; thus the quotient
is a literal polynomial. Mellin transformation gives
\[
 p_h(D)s_hu-s_hP(T_h)u
       =\Theta\bigl(K_{P,h}(u;D)\phi_*\bigr).
 \tag{HCA.11}
\]
Both sides have transform $g(s)K_{P,h}(u;s)$, so Mellin injectivity
proves this equality. This is the full original theta boundary,
not only equality of its classes. Tensoring (HCA.11) with the
other source factors gives its boundaries in (TO.3) with the
original tensor differential and its Koszul signs.

The Mellin isomorphism of (TO.7), (CAU.9) sends $\sigma_hu$
on its offcritical part to precisely its prescribed local jets
$u_\rho$. Equations (HCA.7)--(HCA.8) therefore give the following
endomorphism of the already constructed offcritical direct sum:
\[
 \bigoplus_{\rho\in U:g(\rho)=0}
       \mathcal O_{\mathbb{C},\rho}/(g)
 \longrightarrow
 \bigoplus_{\rho\in U:g(\rho)=0}
       \mathcal O_{\mathbb{C},\rho}/(g),\qquad
 (u_\rho)_\rho\longmapsto
       (-(s-1/2)^{-2}u_\rho)_\rho.
 \tag{HCA.12}
\]
Each denominator is a local analytic unit. Finite support is
unchanged. Thus this is exactly the same reciprocal action as
the one defined in the ambient finite images in $Q$.

\subsection{The full half-trace and the exact global moment limit}

For the packet $Z$, let $\mathcal B_Z$ be its distinct reciprocal
nodes. Each fibre of $\rho\mapsto\beta_\rho$ is the pair
$\rho,1-\rho$, since equality of nonzero reciprocal squares is
equality up to sign. Reflection gives the same full order $m_\beta$
at both points. Define
\[
 f_Z(X)=\prod_{\beta\in\mathcal B_Z}(X-\beta)^{m_\beta},
 \qquad\mathcal A_Z=\mathbb{C}[X]/(f_Z).
 \tag{HCA.13}
\]
The inverse local map (HCA.8) proves that the kernel of
$P\mapsto P(T_h)$ is exactly $(f_Z)$. On choosing one point
in each reflected pair it also gives the full module isomorphism
$E_h\simeq\mathcal A_Z\oplus\mathcal A_Z$. A different choice
interchanges the corresponding two factors; both remain present.
Triangular multiplication in each local basis gives, for $\nu=0,1$,
\[
 \mathcal L_{Z,\nu}(PQ)
 :=\tfrac12\operatorname{Tr}_{E_h}
          \bigl(T_h^{\nu+1}P(T_h)Q(T_h)\bigr)
 =\sum_{\beta\in\mathcal B_Z}m_\beta\beta^{\nu+1}P(\beta)Q(\beta).
 \tag{HCA.14}
\]
Indeed every diagonal entry on a length-$m_\beta$ block equals
$\beta^{\nu+1}P(\beta)Q(\beta)$, and there are two reflected
blocks. This proves the factor $1/2$ and the retained multiplicity.
The radical of this bilinear form on $\mathcal A_Z$ is exactly
its nilradical $(\operatorname{rad}f_Z)/(f_Z)$: vanishing at
every node is sufficient by the sum, and Lagrange interpolation
at a node where a value is nonzero proves necessity, since
$m_\beta\beta^{\nu+1}\ne0$. Thus (HCA.8) retains all local
jets before the trace and (HCA.14) specifies its exact scalar kernel.

Packet growth has a specific trace law. From (HCA.9) and the
idempotent splitting of the larger packet,
\[
 \begin{gathered}
 \tfrac12\operatorname{Tr}_{E_H}
  \bigl(M_{e_{hH}}T_H^{\nu+1}P(T_H)Q(T_H)\bigr)
                  =\mathcal L_{Z,\nu}(PQ),\\
 \mathcal L_{Z',\nu}(PQ)=\mathcal L_{Z,\nu}(PQ)
                  +\mathcal L_{Z'\setminus Z,\nu}(PQ)
                      \qquad(Z\subseteq Z').
 \end{gathered}
 \tag{HCA.15}
\]
The first is a trace on the old direct summand, conjugate by
$\iota$ and $\pi$ to the original endomorphism. The second is
additivity of trace on the two complementary full summands.
Thus uncompressed traces on growing packets are not asserted
constant. Their precise change is the added original packet.

The growth calculation (WBR9) and the complete product proof
(GF1)--(GF13) for this same $F$ give
distinct nonzero reciprocal nodes $\beta_j$, their full orders $m_j$,
and
\[
 \sum_jm_j|\beta_j|<\infty,\quad
 F(w)=a_0\prod_j(1-\beta_jw)^{m_j},\quad
 b_n=\sum_jm_j\beta_j^{n+1}.
 \tag{HCA.16}
\]
That proof applies the explicit bound
$\log M_F(R)=O(\sqrt R\log R)$ with exponent $3/4<1$.
Its circle-mean calculation bounds the zero count by $CR^{3/4}$;
dyadic annuli then prove reciprocal summability. The uniformly
convergent product has exactly the original zeros. Dividing $F$
by it gives a nonvanishing entire function $e^H$. The selected
circles in (GF7)--(GF9) bound $\Re H$ by $CR^{3/4}\log R$;
the circle Fourier formula (GF10)--(GF11) forces every positive
Taylor coefficient of $H$ to vanish. Evaluation at zero fixes
the remaining constant to $a_0$. Thus no unproved factorization
input, exponential factor, or change of mass enters (HCA.16).
The complete product proof is included with the attachment. The nodes are
bounded, since absolute summability excludes infinitely many
nodes outside any disc about zero. Hence, for any fixed $P,Q$,
the sum in (HCA.14) over all nodes is absolutely convergent,
also for $\nu=1$. Exhausting by finite full reflection- and
conjugation-stable packets proves the exact limit
\[
 \begin{gathered}
 \mathcal L_\nu(PQ)=\lim_Z\mathcal L_{Z,\nu}(PQ)
     =\sum_jm_j\beta_j^{\nu+1}P(\beta_j)Q(\beta_j),\\
 \mathcal L_\nu(X^rX^s)=b_{r+s+\nu},\qquad
 \mathcal Q_\nu(P):=\mathcal L_\nu(P^2)
            =c^{\mathsf T}H_{15}^{(\nu)}c
                 \quad\left(P(X)=\sum_{r=0}^{15}c_rX^r\right).
 \end{gathered}
 \tag{HCA.17}
\]
Absolute convergence proves independence of the exhaustion and
justifies the finite polynomial expansion. For real coefficients
the value is real by conjugation of the full node set. The square
in this identity is the algebraic square $P(\beta)^2$.

Equations (HCA.3)--(HCA.4) and (HCA.5)--(HCA.17) are therefore
two equal evaluations of the same original data: the first through
the theta integral and coefficient recurrence, the second through
the full packet action and absolutely convergent half-trace.
Their equality is proved by (HCA.16), not imposed as a new
relation between independently adjustable matrices.
Adjoining these specified observation vertices and the displayed
maps to the path representation $\mathbf C_U$ of (TO.4) gives
one represented extension: every path is evaluated by composition
of the actual maps, and (HCA.9), (HCA.11), (HCA.15), (HCA.17)
prove the new path relations. The old represented diagram is
retained through its identity inclusion. No quotient identifies
its offcritical component with a scalar or with a zero target.

In particular the global trace splits into the absolutely convergent
critical and offcritical sums. The critical-only projection of
(WBR37)--(WBR40) kills every full offcritical packet but need not
kill its trace: its reciprocal primary idempotent $e_\beta$ has
$\operatorname{Tr}_{\mathcal A_Z}M_{X^{\nu+1}e_\beta}
=m_\beta\beta^{\nu+1}\ne0$. Thus the full functional has this
explicit omitted summand under critical projection. At the split
support level all linear maps above send $(\mathfrak a,u)$ to
$(\mathfrak a,Lu)$ and $\tau$ to $\tau$; a killed class lands
at $(\mathfrak a,0)$. The moment and quadratic evaluations remain
the marked observations in their actual linear or polynomial types.

\subsection{The certified finite restriction on this attached observation}

The certified theta calculation (TC1)--(TC28) encloses exactly the
moments (HCA.2). It retains theta terms $1\leq n<20$ on $0\leq y\leq4$,
and, for $J=0,2,\ldots,64$, adds the proved nonnegative tails
\[
 \begin{gathered}
 \frac{16\pi^2 4^{J+1}e^{18}20^4e^{-400\pi}}
                    {1-e^{4/20-40\pi}},\\
 \frac{16\pi^2 4^Je^{18-\pi e^8}}
 {(1-e^{4-2\pi})(2\pi e^8-9/2-J/4)}.
 \end{gathered}
 \tag{HCA.18}
\]
These are exactly (TC5) at $N=20,L=4$, with $N$ the first omitted
integer. All denominators are strictly positive as proved there.
The finite integrands, quadrature subdivisions and dyadic interval
endpoints are preserved in the certificate. The interval recurrence
(TC11)--(TC13) encloses the exact coefficients in (HCA.4) and
their $LDL^{\mathsf T}$ decompositions in the unchanged coordinate
order. Every one of its thirty-two pivot intervals has a positive
lower endpoint, proving
\[
 H_{15}^{(0)}>0,\qquad H_{15}^{(1)}>0.
 \tag{HCA.19}
\]
The certificate is the actual file
\texttt{certificate\_dimension16\_1024.json}, with SHA-256
\nolinkurl{869ebacd3f82e46c3fe5fac1f07bcbaae7f4c333e1a5ba2e0ffcb1f91884a0fc}.
Its exact endpoint checks and the full analytic tail and validated
quadrature proof accompany this statement.

The exact rational endpoint check also gives a uniform quantitative
margin in the original monomial coefficient frame. Write the exact
decompositions as $H_{15}^{(\nu)}=L_\nu D_\nu L_\nu^{\mathsf T}$.
The sum of the absolute upper endpoint bounds in every off-diagonal
row and column of $L_\nu-I$ is strictly less than $1/64$.
Every diagonal pivot of $D_0$ is greater than $6\cdot10^{-107}$,
and every pivot of $D_1$ is greater than $7\cdot10^{-111}$.
These comparisons use exact dyadic fractions in the certificate.
For a matrix $E$, Cauchy--Schwarz with row weights proves
$\|Ex\|_2^2\leq\|E\|_\infty\|E\|_1\|x\|_2^2$:
bound the square of row $i$ by
$(\sum_j|E_{ij}|)(\sum_j|E_{ij}||x_j|^2)$ and sum in $i$.
Thus $\|L_\nu-I\|_2<1/64$ and
$\|L_\nu^{\mathsf T}x\|_2>(63/64)\|x\|_2$ for $x\ne0$.
Substitution in the exact decomposition proves
\[
 \begin{gathered}
 H_{15}^{(\nu)}\succ\gamma_\nu I_{16},\qquad
 \gamma_0=(63/64)^2\,6\cdot10^{-107},\quad
 \gamma_1=(63/64)^2\,7\cdot10^{-111}.
 \end{gathered}
 \tag{HCA.19a}
\]
No polynomial basis or mass was changed in obtaining this bound.

Consequently, every nonzero real polynomial of degree at most
fifteen satisfies $\mathcal Q_\nu(P)>0$ for both $\nu=0,1$.
Padding its coefficient vector with zeros identifies its form
with the corresponding leading principal submatrix; positivity
of (HCA.19) proves the assertion. More generally, on the explicitly
specified finite cone
\[
 \mathscr C_{15}=\left\{\sum_j A_j(X)^2+X\sum_j B_j(X)^2:
          A_j,B_j\in\mathbb{R}[X]_{\leq15},
          \text{finite sums}\right\},
 \tag{HCA.20}
\]
one has $\mathcal L_0(C)>0$ for every nonzero $C$ in that cone.
Indeed it is the sum of the positive forms in (HCA.19), and a
nonzero polynomial represented in the cone has at least one
nonzero summand. This uses the shifted matrix through
$\mathcal L_0(XB^2)=\mathcal Q_1(B)$; no characterization of
larger polynomial cones or positivity at higher degrees is presumed.

\subsection{The precise exclusion for an offcritical separator}

Take any actual offcritical pair $\beta,\overline\beta$ of
multiplicity $m$ in (HCA.16). In original zero coordinates it is
\[
 \beta=-\frac1{(\rho-1/2)^2}
 =\frac{\gamma^2-\delta^2+2i\gamma\delta}
                       {(\gamma^2+\delta^2)^2},
 \qquad\rho=1/2+\delta+i\gamma.
 \tag{HCA.21}
\]
It is nonreal because $\delta\gamma\ne0$. Put $r=|\beta|/2$,
let $\mathcal A$ be the set of all other distinct nodes of modulus
at least $r$, and define
\[
 B_\beta(X)=\prod_{\alpha\in\mathcal A}(X-\alpha),\quad
 B_0=\prod_{\alpha\in\mathcal A}(r+|\alpha|),\quad a=\#\mathcal A.
 \tag{HCA.22}
\]
The set is finite by (HCA.16), conjugation stable, and excludes
$\beta,\overline\beta$, so $B_\beta$ is real and
$B_\beta(\beta)\ne0$. For $\nu=0,1$ choose either $c_\nu$ with
$\beta^{\nu+1}c_\nu^2=-1$, and retain the actual constants
\[
 C_\nu=\frac{|c_\nu|B_0}{|B_\beta(\beta)|}
        \left(1+\frac{|\Re\beta|+r}{|\Im\beta|}\right),\qquad
 T_\nu=\sum_{|\beta_j|<r}m_j|\beta_j|^{\nu+1}.
 \tag{HCA.23}
\]
Both tails converge and $T_1\leq rT_0$. Define for every integer
$N\geq0$
\[
 \begin{gathered}
 w_{\nu,N}=\frac{c_\nu}{\beta^NB_\beta(\beta)},\quad
 v_{\nu,N}=\frac{\Im w_{\nu,N}}{\Im\beta},\quad
 u_{\nu,N}=\Re w_{\nu,N}-(\Re\beta)v_{\nu,N},\\
 P_{\nu,N}(X)=X^NB_\beta(X)(u_{\nu,N}+v_{\nu,N}X),
 \qquad D_{\nu,N}=\deg P_{\nu,N}\leq N+a+1.
 \end{gathered}
 \tag{HCA.24}
\]
This is precisely (TC24) for $\nu=0$. The shifted version uses
the original shifted form and its exponent $\nu+1=2$.
The polynomial is real and nonzero, because its value at $\beta$
is $c_\nu\ne0$. Its actual degree is $N+a+1$ when
$v_{\nu,N}\ne0$, and $N+a$ otherwise. Evaluation proves
$P_{\nu,N}(\beta)=c_\nu$ and conjugate equality at
$\overline\beta$. Hence its selected two contributions to
the actual scalar form are exactly $-2m$.

Every other node of modulus at least $r$ is a zero of the
polynomial. For $|X|\leq r$, the elementary estimates
$|B_\beta(X)|\leq B_0$ and
$|u_{\nu,N}+v_{\nu,N}X|\leq|w_{\nu,N}|+
(|\Re\beta|+r)|w_{\nu,N}|/|\Im\beta|$ give
$|P_{\nu,N}(X)|\leq C_\nu2^{-N}$. Absolute summation proves
the exact identity and tail bound
\[
 \begin{gathered}
 \mathcal Q_\nu(P_{\nu,N})=-2m+\mathcal T_{\nu,N},\qquad
 \mathcal T_{\nu,N}=\sum_{|\beta_j|<r}
                  m_j\beta_j^{\nu+1}P_{\nu,N}(\beta_j)^2\in\mathbb{R},\\
 |\mathcal T_{\nu,N}|\leq C_\nu^2T_\nu4^{-N}.
 \end{gathered}
 \tag{HCA.25}
\]
No tail is assigned a sign before applying the certificate.
If $D_{\nu,N}\leq15$, the nonzero polynomial in (HCA.24) is
within (HCA.19), and therefore
\[
 \boxed{\mathcal T_{\nu,N}=2m+\mathcal Q_\nu(P_{\nu,N})>2m,
 \qquad C_\nu^2T_\nu4^{-N}>2m.}
 \tag{HCA.26}
\]
This is the finite exclusion on the actual packet together with its
entire reciprocal tail. In particular, when $a\leq14$, choosing
$N=14-a$ makes its degree at most fifteen and gives
\[
 C_\nu^2T_\nu>2m\,4^{14-a}\qquad(\nu=0,1).
 \tag{HCA.27}
\]
If the affine factor has zero leading coefficient, the exact
degree condition in (HCA.26) additionally retains that case.
There is also an explicit strictly positive margin. For a polynomial
$P(X)=\sum_{j=0}^D p_jX^j$ with $D\leq15$,
Cauchy--Schwarz gives
$|P(\beta)|^2\leq(\sum_j|p_j|^2)(\sum_{j=0}^D|\beta|^{2j})$.
For $P=P_{\nu,N}$ its value is $c_\nu$ and
$|c_\nu|^2=|\beta|^{-\nu-1}$. Applying (HCA.19a) to its
coefficient vector padded with zeros strengthens (HCA.26) to
\[
 \boxed{C_\nu^2T_\nu4^{-N}\geq\mathcal T_{\nu,N}
 >2m+\frac{\gamma_\nu|\beta|^{-\nu-1}}
               {\sum_{j=0}^{D_{\nu,N}}|\beta|^{2j}}
 \quad(D_{\nu,N}\leq15).}
 \tag{HCA.27a}
\]
The inequality on the left uses the actual tail's strict positivity
already proved in (HCA.26), followed by its absolute upper bound.

The shifted certificate also tests the identical original polynomial
$P_{0,N}$, without changing $c_0$ or any coefficient. Its two selected
terms in $\mathcal Q_1$ are
$m\beta^2c_0^2+m\overline\beta^2\overline{c_0}^{\,2}
=-2m\Re\beta$. The absolute value of its remaining sum is at
most $C_0^2T_14^{-N}\leq rC_0^2T_04^{-N}$. If
$D_{0,N}\leq15$ and $p$ is its original coefficient vector padded
with zeros through degree fifteen, (HCA.19a) consequently proves
the simultaneous exact bounds
\[
 \begin{aligned}
 C_0^2T_04^{-N}
 &>\max\left\{2m+\gamma_0\|p\|_2^2,
          \frac{2m\Re\beta+\gamma_1\|p\|_2^2}{r}\right\}\\
 &>\max\left\{2m,\frac{4m\Re\beta}{|\beta|}\right\}.
 \end{aligned}
 \tag{HCA.27b}
\]
The second entry remains valid even when its numerator is negative:
the absolute upper tail bound is still an upper bound for the signed
tail. Both positive coefficient margins are strict. Substitution
$\|p\|_2^2\geq|\beta|^{-1}/\sum_{j=0}^{D_{0,N}}|\beta|^{2j}$
from the preceding evaluation argument retains an entirely explicit
margin in both entries of the first maximum. The elementary shifted
entry improves $2m$ exactly when $\Re\beta/|\beta|>1/2$.
In the original coordinates (HCA.21) that inequality is
$\gamma^2>3\delta^2$, since
$\Re\beta/|\beta|=(\gamma^2-\delta^2)/(\gamma^2+\delta^2)$.
This additional comparison uses the same original unshifted
separator, as well as the separately constructed shifted family.

There is a further comparison between the two separately phased
families. Set $U_{\nu,N}=C_\nu^2T_\nu4^{-N}$. The exact modulus
$|c_\nu|^2=|\beta|^{-\nu-1}$ in (HCA.23) gives
$C_1^2=C_0^2/|\beta|$. Since $T_1\leq rT_0$ and
$r=|\beta|/2$, one obtains $U_{1,N}\leq U_{0,N}/2$.
For $D_{1,N}\leq15$, apply (HCA.19a) to the original coefficient
vector $p_1$ of $P_{1,N}$, padded through degree fifteen, and use
(HCA.25). The resulting exact inequalities are
\[
 \begin{gathered}
 U_{0,N}\geq2U_{1,N}
 >4m+2\gamma_1\|p_1\|_2^2
 \geq4m+
       \frac{2\gamma_1|\beta|^{-2}}
            {\sum_{j=0}^{D_{1,N}}|\beta|^{2j}},\\
 a\leq14\quad\Longrightarrow\quad
 C_0^2T_0>4m\,4^{14-a}=m\,4^{15-a}.
 \end{gathered}
 \tag{HCA.27c}
\]
Here the last line chooses $N=14-a$, so that $D_{1,N}\leq15$;
the coefficient estimate is the same evaluation argument as in
(HCA.27a). This estimate uses the shifted polynomial $P_{1,N}$
to bound the original unshifted tail constant. It is distinct from
the same-polynomial comparison (HCA.27b), which retains $P_{0,N}$.

For the original unshifted quantitative separator of (TC23), put
\[
 N_*=
 \begin{cases}0,&T_0=0,\\
 \displaystyle\max\!\left\{0,1+
  \left\lfloor\frac{\log(C_0^2T_0/m)}{\log4}\right\rfloor\right\},
       &T_0>0.
 \end{cases}
 \tag{HCA.28}
\]
Its definition gives $C_0^2T_04^{-N_*}\leq m$ and thus
$\mathcal Q_0(P_{0,N_*})\leq-m<0$ by (HCA.25). Comparison
with (HCA.19) proves the concrete restriction
\[
 \boxed{\deg P_{0,N_*}\geq16,\qquad N_*+a+1\geq16.}
 \tag{HCA.29}
\]
The shifted comparison strengthens the nominal bound for the same
original threshold. If $a\leq14$, (HCA.27c) implies $T_0>0$ and
$\log(C_0^2T_0/m)/\log4>15-a$. Because $15-a$ is an integer,
its floor is at least $15-a$, and (HCA.28) gives
\[
 \boxed{a\leq14\quad\Longrightarrow\quad
        N_*\geq16-a,\qquad N_*+a+1\geq17.}
 \tag{HCA.29a}
\]
This is a bound on the nominal degree. Its actual degree remains
$N_*+a+1$ if $v_{0,N_*}\ne0$ and $N_*+a$ otherwise; the latter
case can still have degree sixteen. There is a further exact boundary
case: if $a\leq15$ and $v_{1,15-a}=0$, the shifted polynomial with
$N=15-a$ has degree exactly fifteen. The first line of (HCA.27c)
then gives $C_0^2T_0>m4^{16-a}$, and the identical integer-floor
argument yields $N_*\geq17-a$. No assertion on a matrix of degree
sixteen is needed for either conclusion.

The rational rounding of (TC25) keeps the actual quadratic form
strictly negative and has degree at most $N_*+a+1$. Its actual
degree is therefore at least sixteen as well. Coefficient rounding
cannot produce a smaller-degree negative witness against the
certified matrices. The same conclusion holds for any rational
negative polynomial in either shifted family, by (HCA.19).

The full finite packet alone may already contribute $-2m$ under
the exact half-trace (HCA.14); its remaining original arithmetic
nodes still enter (HCA.17). The certificate therefore excludes
the low-degree global separators specified above, rather than
deleting that offcritical packet from $T_UQ$. It imposes no
positivity assertion on $H_{16}$ or any larger matrix. Every
source block, its full reciprocal coordinate, its scalar radical,
its receiving supported zero and its exact ambient trace remain
the explicitly constructed parts of this attachment.
